\section{CoughKD-ShiftAudit}
\label{sec:method}

\subsection{Problem Setup}

Let $\mathcal{D}_s=\{(x_i,y_i,u_i)\}$ be a source cough dataset with recording $x_i$, label $y_i$, and subject identifier $u_i$. Let $\mathcal{D}_t$ be a target dataset used only for external evaluation and stress testing. The goal is not to maximize a single in-domain score, but to audit whether a compact student $f_s$ distilled from a high-capacity teacher $f_t$ remains reliable under $\mathcal{D}_s \rightarrow \mathcal{D}_t$ shift.

We treat the source split as subject-disjoint. Target labels are used for evaluation only, not for target-supervised model selection. When a target contains multiple clips from one subject or original recording, we report both clip-level and subject-level views when possible.

\subsection{Teacher and Student}

The current teacher is PANNs CNN14 at 16 kHz, initialized from AudioSet-pretrained weights and adapted to the source label space \cite{kong2020panns}. The student is a compact depthwise CNN with 20.7K parameters. This gives an extreme compression ratio of about 3912x relative to the 81M-parameter teacher. The teacher/student pair is intentionally simple: the paper audits whether even standard KD assumptions are reliable before expanding to a larger teacher-by-student grid.

\subsection{Audited Training Variants}

We compare a bounded set of training variants:
\begin{itemize}
    \item CE-only student training;
    \item vanilla response KD;
    \item source-only continuation;
    \item target-consistency KD;
    \item confidence-gated target-consistency KD;
    \item shortcut-suppressed KD;
    \item disagreement-gated KD;
    \item probe-adversarial KD.
\end{itemize}

These variants are not presented as a new SOTA architecture search. They form an audit set for common KD assumptions: that teacher probabilities help, that target consistency helps, that shortcut suppression helps, and that domain readability reduction helps.

\subsection{Audit Evidence}

\textsc{CoughKD-ShiftAudit} reports:
\begin{itemize}
    \item external macro AUROC, COVID AUROC, and macro AUPRC;
    \item ECE, Brier score, and sensitivity at 95\% specificity;
    \item domain probe AUC and task probe AUC over student representations;
    \item parameter count, checkpoint size, and latency;
    \item metadata-defined COUGHVID slices;
    \item bootstrap target-subset stability;
    \item target-unlabeled guard stress tests;
    \item clip-vs-subject aggregation on segmented targets;
    \item source-target overlap control for external and auxiliary manifests.
\end{itemize}

\subsection{Overlap Control}

External validation claims are invalid if the target overlaps with the source. Before reporting any new external or auxiliary target, we run a manifest overlap audit against Coswara. The audit checks exact recording IDs, subject IDs, paths, path stems, path tails, shared identifier-like tokens, and optionally SHA-256 audio hashes. DiCOVA is treated as a Coswara-related auxiliary protocol unless overlap is explicitly controlled.
